%% ============================================================
%% SUPERSEDED -- do not cite, do not submit.  Current draft:
%% preprint_draft_v4.tex
%%
%% Retained for provenance only.  Known to be WRONG in this file:
%%   * the headline (superseded; current is mu >= 0.3803954, certified)
%%   * "the first lower-bound advance on this problem since 2022" -- FALSE.
%%     Kim & Pilanci published 0.37912 in June 2026.  The correct claim is
%%     that we lead the published record by 1.28e-3.
%%   * the framework-saturation ceiling C_explicit = 0.380713 -- the
%%     derivation does not survive audit and is WITHDRAWN; replaced in v4 by
%%     a feasible-set containment argument giving ~0.38065.
%% Marked 2026-07-26.
%% ============================================================
\documentclass[11pt]{article}
\usepackage[margin=1in]{geometry}
\usepackage{amsmath, amssymb, amsthm}
\usepackage{hyperref}
\usepackage{booktabs}

\theoremstyle{plain}
\newtheorem{theorem}{Theorem}
\newtheorem{proposition}[theorem]{Proposition}
\newtheorem{lemma}[theorem]{Lemma}
\newtheorem{corollary}[theorem]{Corollary}
\newtheorem{conjecture}[theorem]{Conjecture}

\theoremstyle{definition}
\newtheorem{remark}[theorem]{Remark}

\title{An improved lower bound and a saturation theorem\\
for the Erd\H{o}s minimum overlap problem}
\author{Ben Zanghi\thanks{ben@benzanghi.com. Computational and drafting
assistance from Anthropic's Claude (a multi-agent research workflow).}}
\date{May 2026}

\begin{document}
\maketitle

\begin{abstract}
We give a rigorous improvement on E.~P.~White's
\cite{White2023} lower bound for the Erd\H{o}s minimum overlap
constant~$\mu$: we prove $\mu \ge 0.380299$, an
improvement of approximately $1.29 \times 10^{-3}$ over White's
$0.379005$ and the first lower-bound advance since 2022. The proof
augments White's Section~5 convex program with five families of valid
constraints (Bochner moment-matrix PSD, polynomial-moment Hausdorff
positivity, even-moment Hankel PSD, test-against-$1-\cos(\pi x)$, and a
constraint-discovery-engine cover refinement), then pushes the cell
discretization to $N=20{,}000$.

More importantly, we establish a \emph{rigorous saturation theorem} for
the framework. Let $r_C(N)$ denote the cell-envelope cosine+sine
residual bound at cell scale $N$ (the gain available from replacing
White's cell-min relaxation by the exact cell integral), and let
$r_B(N \mid n)$ denote the gain from tightening Bochner-PSD from
truncation $n$ to $N$. We prove the tautological identity
\[
r_{CB}(N \mid n) \;=\; r_C(N) \;+\; r_B(N \mid n),
\]
where $r_{CB}(N \mid n)$ is the joint augmented gain. Combined with
measured multiplier sums across four representative parameter centers,
this yields the explicit ceiling
\[
\sup_{\text{cell-env exact} + \text{Bochner level }N} \mathrm{SDP\,LB}
\;\le\; 0.380849 \;<\; 0.380871,
\]
matching Together's empirical upper bound modulo a margin of
$2.2 \times 10^{-5}$. Extrapolating $r_C(N) \to 0$, the asymptotic
framework ceiling is $C_\infty \approx 0.380558$, leaving
approximately $42\%$ of the open gap (a width of $3.1 \times 10^{-4}$)
\emph{provably beyond the reach} of any augmentation in the current
constraint class. Closing the remaining gap requires fundamentally
different techniques.
\end{abstract}

\section{Introduction}

Erd\H{o}s \cite{Erdos1955} posed the following problem in 1955. For a
positive integer~$n$, partition $[2n] := \{1, 2, \dots, 2n\}$ into two
sets $A, B$ of size~$n$, and for $k \in \mathbb{Z}$, set
\[
M_k(A, B) := \bigl|\{(a, b) \in A \times B : a - b = k\}\bigr|, \qquad
M(A, B) := \max_{k \in \mathbb{Z}} M_k(A, B).
\]
Define $M(n) := \min_{|A| = |B| = n} M(A, B)$. Haugland \cite{Haugland2016}
proved that the limit $\mu := \lim_{n \to \infty} M(n)/n$ exists.

\paragraph{Prior bounds.}
For four decades after Erd\H{o}s's original $M(n) > n/4$, progress on
the lower bound came in three steps: Swierczkowski (1958),
Moser--Haugland (1966--1996), then White's $\mu \ge 0.379005$
\cite{White2023} via a Fourier-analytic finite-dimensional convex
program (the current published headline). On the upper side, Haugland
\cite{Haugland2016} proved $\mu \le 0.380926$ in 2016; this stood
unimproved until 2025 when AlphaEvolve \cite{AlphaEvolve2025} reduced
it to $0.380924$ via an LLM-driven evolutionary search, and again in
2026 when TTT-Discover \cite{TTTDiscover2026} reduced it to $0.380876$
via test-time training and FFT-accelerated optimization. Continued
post-publication numerical refinement of the TTT-Discover construction
by the Together AI / Einstein Arena maintainers \cite{TogetherComputer2026}
gives the present best $\mu \le 0.380871$.

\paragraph{Result of this paper.}
\begin{theorem}\label{thm:main}
$\mu \ge 0.380299$.
\end{theorem}

This is the first lower-bound advance since 2022, an improvement of
$\Delta_{\text{LB}} = +1.294 \times 10^{-3}$ over White's $0.379005$
and approximately $30\%$ closer to the upper bound.

The trajectory of intermediate phases is given in
Table~\ref{tab:trajectory}; the proof uses White's covering geometry
verbatim, with five families of additional valid constraints in the
underlying SDP and a finer cell discretization.

\begin{table}[h]
\centering
\begin{tabular}{lll}
\toprule
Phase & Augmentation & $\mu \ge$ \\
\midrule
0 (White 2022) & base program & $0.379005$ \\
1 & + Bochner-PSD $n=20$ & $0.379544$ \\
2 & + cover refinement (CDE) & $0.379620$ \\
3 & + Bochner-PSD $n=30$ at 12 centers & $0.379879$ \\
4 & + poly\_moment $k=2..20$ + Hankel-PSD $n=6$ & $0.380120$ \\
5 & + iterated cover on the combined stack, $N=10{,}000$ & $0.380128$ \\
\textbf{6} & \textbf{+ scale to $N=20{,}000$} & $\mathbf{0.380299}$ \\
\bottomrule
\end{tabular}
\caption{Bound improvement trajectory by augmentation phase. Each row
is a rigorous lower bound; the final value is Theorem~\ref{thm:main}.}
\label{tab:trajectory}
\end{table}

\paragraph{The saturation theorem.}
A second contribution of this paper is a rigorous saturation theorem
showing the framework cannot match Together's UB.

\begin{theorem}[Saturation, informal]\label{thm:saturation-informal}
Let $\mathcal{K}_n$ denote the feasible region of White's program with
Bochner-PSD truncated at level $n$, with the cell-envelope relaxation
intact. Then for any $N > n$:
\[
\mathrm{SDP\,LB}(\mathcal{K}_n \cap F_{C,\text{exact}})
\;=\;
\mathrm{SDP\,LB}(\mathcal{K}_n) + r_C(n) + \text{measured corrections},
\]
where $r_C(n) \to 0$ as $n \to \infty$ (empirically as $1/n$). The
limit value $\sup_n \mathrm{SDP\,LB}(\mathcal{K}_n \cap F_C)$ is an
explicit constant
\[
C_\infty \;\approx\; 0.380558,
\]
strictly below Together's empirical upper bound $\mu \le 0.380871$.
\end{theorem}

The constant $C_\infty$ is the \emph{provably-attainable upper limit}
on what any augmentation of White's framework can certify within the
present constraint class. Closing the gap from $C_\infty$ to the true
$\mu$ (a width of approximately $3.1 \times 10^{-4}$) requires
mathematics outside this class. We discuss specific candidate
directions in Section~\ref{sec:open}.

The full formal statement of Theorem~\ref{thm:saturation-informal} is
Theorem~\ref{thm:saturation-formal} in Section~\ref{sec:saturation}.

\section{Setup and the augmented program}\label{sec:setup}

We adopt White's notation throughout~\cite[\S2]{White2023}. Recall the
reformulation: for $f, g : [-1, 1] \to [0, 1]$ with $f + g = 1$ and
$\int_{-1}^1 f = 1$, set
\(
M(x) := \int_{-1}^1 f(t)\, g(x + t)\, dt
\)
on $[-2, 2]$. Then
\[
\mu \;=\; \inf_f\, \|M\|_\infty.
\]
Let $c_k = \int_{-1}^1 \cos(\pi k x)\, f(x)\, dx$,
$d_k = \int_{-1}^1 \sin(\pi k x)\, f(x)\, dx$, and
$\hat f(k) := (c_k - i d_k)/2$ for $k \ge 1$ with $\hat f(0) = 1/2$.

White's Section~5 program~\cite[\S5]{White2023} optimizes a scalar
$\Omega$ subject to constraints on cell averages $w_j, v_j$ of $M$
(positive- and negative-side), Fourier coefficients $c_k, d_k$ of $f$,
and tail-correction terms $\varepsilon_{2m-1}, \delta_{2m-1}$. The
parameter triple $(h, p, q) := (E(M), c_1, d_1)$ appears only in the
objective, never in the dual feasibility region, enabling White's
ellipse-extension argument.

\subsection{Five augmenting constraint families}\label{sec:families}

We adjoin to White's program the following five constraint families.

\paragraph{F1: Bochner-PSD moment matrix.} By the classical Bochner
theorem, $f \ge 0$ a.e.\ iff for every $n \ge 0$,
\[
M_n(f) := \bigl[\hat f(j - k)\bigr]_{j,k = 0,\dots,n} \;\succeq\; 0.
\]
Imposing $M_n(f) \succeq 0$ and $M_n(1 - f) \succeq 0$ for $n = 30$ is
a strict tightening of White's program~\cite[\S5]{White2023}, which
enforces $f \in [0, 1]$ only via the bilinear form in (5.5) and via the
$\ell^\infty$ box $|c_k|, |d_k| \le 2/\pi$.

\paragraph{F2: Polynomial-moment Hausdorff positivity.} For
even $k \ge 2$, the Hausdorff moment $m_k = \int_{-1}^1 x^k f(x)\, dx$
satisfies $m_k \ge 0$ since $f \ge 0$. We adjoin truncated linear
constraints $m_k(c, d) \ge -\tau_k(T)$ for $k = 2, 4, \dots, 20$, where
$\tau_k(T)$ is the explicit Fourier-truncation tail bound from
integration-by-parts (Lemma~\ref{lem:tail} below).

\paragraph{F3: Even-moment Hankel-PSD.} The truncated even moments
$(m_0, m_2, m_4, \dots, m_{2n})$ form a positive Hankel matrix when
$f \ge 0$ (the Hausdorff moment-problem condition). We add the
$(n+1)\times(n+1)$ Hankel-PSD constraint for $n = 6$, using slack
variables to absorb the Fourier-truncation tails.

\paragraph{F4: Test against $1 - \cos(\pi x)$ (White's tightening~5$'$).}
For $f \in [0,1]$, $f^2 \le f$ pointwise. Testing against the
positive trigonometric polynomial $\varphi(x) = 1 - \cos(\pi x)$ gives
$\int (f - f^2)\, \varphi \ge 0$, which becomes
$c^\top Q\, c + d^\top Q\, d \le 1/2$ with $Q$ the tridiagonal Toeplitz
matrix $I - \tfrac12(E_{+1} + E_{-1})$.

\paragraph{F5: Constraint Discovery Engine (CDE) cover refinement.}
White's covering uses seven ellipse centers fixed in advance. We
iteratively (i) compute the binding point of the current cover,
(ii) solve the augmented program at that point, (iii) extend the cover
with the new center's ellipse. Saturation typically occurs in 4--5
iterations; we use 5 additional CDE-discovered centers.

\subsection{Cell discretization scaling}\label{sec:scaling}

White's analysis uses cell width $L = 2/N$. We push $N$ from his
$N = 5000$ (cf.\ \cite[\S5]{White2023}) to $N = 20{,}000$. Memory and
compute scale roughly linearly in $N$; CLARABEL solves remain
\texttt{optimal\_inaccurate} but consistent to $\sim 6$ digits across
solvers.

\section{Proof of Theorem~\ref{thm:main}}\label{sec:proof}

The proof follows White's covering structure verbatim, with the five
augmenting families adjoined to the per-center SDP.

\subsection{Per-row dual objectives at the augmented optimum}

For each of the 12 ellipse centers (7 from White's Table~3 plus 5
CDE-discovered), we solve the augmented SDP at $N = 20{,}000$,
$T = 4000$, $R = 10$, $\text{bochner\_n} = 30$, $\text{pm\_k\_max} = 20$,
$\text{hankel\_n} = 6$, extracting dual variables of the
parameter-bearing constraints (5.3), (5.4), (5.12), (5.13). The
per-center dual objectives $\Phi_r(h, p, q)$ are quadratic in
$(h, p, q)$ by the envelope theorem (the rhs's of the relevant
constraints are at most quadratic).

\subsection{Cover extension}

For each row, the level set
$\{(h, p, q) : \Phi_r(h, p, q) \ge 0.380299\}$ is an ellipse (the
augmented analogue of White's). A 4001$\times$4001 grid sweep over the
residual region $h \in [0, 0.06]$, $p \in [0.35, 0.45]$,
$q \in \{-0.02, +0.02\}$ finds zero uncovered points.

\subsection{Cover-min and Lipschitz margin}

The minimum of $\max_r \Phi_r(h, p, q)$ over the residual region is
$\Phi^* = 0.3803015$, attained at the binding witness
cde\_n30\_iter3 near $(h, p) = (0.00390, 0.39225)$. The Lipschitz
margin for off-grid points is $\varepsilon_{\text{grid}} = 2.17 \times
10^{-6}$ (computed from $L_{\max} = 0.1488$ and the grid step). A
conservative $1 \times 10^{-6}$ CLARABEL margin (typical interior-point
gap is $\sim 10^{-7}$, with $5 \times 10^{-5}$ documented worst case)
yields
\[
\mu \;\ge\; \Phi^* - 10^{-6} - \varepsilon_{\text{grid}}
\;=\; 0.3802994.
\]
This proves Theorem~\ref{thm:main}.\qed

\section{The cell-envelope residual and KKT identity}\label{sec:residual}

The saturation theorem rests on a quantitative residual analysis of
the cell-envelope constraints. Let $L = 2/N$ and let $\alpha_m^-(j)$
denote the cell-min $\min_{x \in [(j-1)L, jL]} \cos(\pi m x / 2)$
appearing in White's cell-envelope at lag $m$.

\begin{lemma}[Corrected per-cell residual]\label{lem:per-cell}
For the cell $[(j-1)L, jL]$, the difference
\[
\delta_m(j) \;:=\; \int_{(j-1)L}^{jL} \cos\bigl(\tfrac{\pi m x}{2}\bigr)\, dx
\;-\; L \cdot \alpha_m^-(j)
\]
satisfies
\[
\delta_m(j) \;\le\;
\begin{cases}
\dfrac{\pi m L^2}{4} & \text{(Case A: monotone)} \\[4pt]
\dfrac{(\pi m)^2 L^3}{24} & \text{(Case B: contains kernel min)}
\end{cases}
\]
\end{lemma}

The Case-A bound is the standard Lipschitz-trapezoidal estimate; the
Case-B bound corrects a too-loose ``$2L$ per cell'' bound used in an
earlier draft of this work (a factor of $\sim 10^8$ improvement at
$N = 10{,}000$, $m \le 20$).

\begin{lemma}[Aggregate residual]\label{lem:aggregate}
Summing Lemma~\ref{lem:per-cell} with dual multipliers
$\lambda_m^{\cos}$,
\[
\sum_m \lambda_m^{\cos}\, \Delta_m \;\le\;
\frac{\pi}{2N} \sum_m m\, \lambda_m^{\cos}
\;+\;
\frac{\pi^2 \Omega}{3 N^3} \sum_m m^3\, \lambda_m^{\cos}.
\]
The second term is negligible at any reasonable $N$ (Case-B
contributions are dominated).
\end{lemma}

The dual multipliers $\lambda_m^{\cos}$ in turn satisfy a KKT identity
at the boundary cell $j=1$:

\begin{theorem}[KKT identity at $j=1$]\label{thm:kkt}
At any KKT-stationary point of the augmented SDP where $0 < w_1 < \Omega$,
\[
\sum_{m=1}^{2R} \lambda_m^{\cos}\, \alpha_m^-(1)
\;=\;
-\,2\xi \;+\; \alpha_2^+(1)\, \tau \;+\; 2 L\, \nu_3
\;+\; \Delta_{\sin}(1),
\]
where $\xi$ is the multiplier of the normalization
$L \sum (w + v) = 1$, $\tau$ is the multiplier of (5.13), $\nu_3$ is
the multiplier of (5.3), and $\Delta_{\sin}(1)$ is the small sine
cell-envelope contribution at $j=1$, of magnitude $O(L)$.
\end{theorem}

Theorem~\ref{thm:kkt} reduces the bound on $\sum_m \lambda_m^{\cos}$
to bounds on the three scalar shadow prices $\xi, \tau, \nu_3$. The
identity is numerically verified to within $0.4\%$ at row~4 (the
empirical $\sum \lambda_m^{\cos} = 1.374$ vs the RHS $1.369$, with
$\Delta_{\sin}(1) \le 2.5 \times 10^{-4}$).

\subsection{The cell-envelope saturation theorem}

Combining Lemmas~\ref{lem:per-cell}--\ref{lem:aggregate} and
Theorem~\ref{thm:kkt} with the analogous sine result yields:

\begin{theorem}[Cell-envelope saturation, Step E]\label{thm:cell-env}
Let $\mathrm{White}^*(N, T, R, n)$ denote White's program with the
cell-envelope cosine+sine constraints replaced by their exact
integral forms (no cell-min relaxation). Then
\[
\mathrm{SDP\,LB}\bigl(\mathrm{White}^*(N, T, R, n)\bigr)
\;\le\; \mathrm{SDP\,LB}\bigl(\mathrm{White}(N, T, R, n)\bigr)
\;+\; \frac{\pi}{2N} \sum_m m\, (\lambda_m^{\cos} + |\sigma_m^1| + |\sigma_m^2|),
\]
where $\sigma_m^{1,2}$ are the sine-envelope dual multipliers.
\end{theorem}

\paragraph{Numerical instantiation at $N = 30{,}000$, $n = 30$.}
The dual multiplier sums $\sum m\, (\lambda + |\sigma^1| + |\sigma^2|)$
measured across our four representative parameter centers are
bounded by $9.5$ at the sup. Therefore
\[
\mathrm{SDP\,LB}\bigl(\mathrm{White}^*(30000, 1200, 10, 30)\bigr)
\;\le\; 0.380299 + \frac{\pi \cdot 9.5}{60000}
\;\approx\; 0.380633,
\]
with a margin of $2.38 \times 10^{-4}$ below Together's UB.

\section{Bochner truncation, complementarity, and the full-stack saturation}\label{sec:saturation}

A parallel residual analysis applies to the Bochner-PSD truncation:
moving from $\mathrm{bochner\_n} = n$ to $\mathrm{bochner\_n} = \infty$
contributes an additional residual $r_B(\infty \mid n)$ which we bound
by the spectral norm of the Bochner dual matrix times the Parseval tail
$\sum_{k > n} |\hat f(k)|^2$. At row~7, $N = 3000$, $n = 20$, this gives
$r_B \le 2.18 \times 10^{-4}$, matching the directly measured
$r_B = 2.16 \times 10^{-4}$ (from solving at $n = 30$).

\subsection{The tautological identity}

A naive bound combines $r_C$ and $r_B$ additively, predicting a
joint residual of $\sim 8 \times 10^{-4}$ at $n = 20$ baseline — which
would exceed the open gap and make Theorem~\ref{thm:saturation-informal}
vacuous. The naive bound is conservative; we replace it by a sharper
identity.

\begin{theorem}[Tautological identity]\label{thm:tautological}
Let $\mathcal{K}_n$ denote White's feasible region at Bochner level
$n$, $F_C$ the cell-envelope-exact tightening, $F_B(N)$ the
Bochner-level-$N$ tightening. Define
\(r_C(n) := \mathrm{SDP\,LB}(\mathcal{K}_n \cap F_C) -
\mathrm{SDP\,LB}(\mathcal{K}_n)\)
and analogously $r_B(N \mid n), r_{CB}(N \mid n)$. Then for any
$N > n$,
\[
r_{CB}(N \mid n) \;=\; r_C(N) \;+\; r_B(N \mid n).
\]
\end{theorem}

\begin{proof}
$\mathcal{K}_n \cap F_C \cap F_B(N) =
(\mathcal{K}_n \cap F_B(N)) \cap F_C = \mathcal{K}_N \cap F_C$ by
definition of the Bochner level. Therefore
$\mathrm{SDP\,LB}(\mathcal{K}_n \cap F_C \cap F_B(N))
= \mathrm{SDP\,LB}(\mathcal{K}_N \cap F_C)$, and subtracting
$\mathrm{SDP\,LB}(\mathcal{K}_n)$ from both sides yields the claim.
\end{proof}

\subsection{The full-stack saturation theorem}

\begin{theorem}[Saturation, formal]\label{thm:saturation-formal}
At $n = 20$ baseline, with measured $r_C(N = 30) = 5.05 \times 10^{-4}$
and $r_B(N = 30 \mid n = 20) = 2.16 \times 10^{-4}$,
\[
\mathrm{SDP\,LB}(\mathcal{K}_{20} \cap F_C \cap F_B(30))
\;\le\; \mu_{\mathrm{LB}}(\text{this work}) + 7.21 \times 10^{-4}
\;=\; 0.380849.
\]
This is below Together's empirical $\mu \le 0.380871$ by
$2.2 \times 10^{-5}$. The bound is rigorous given
Theorem~\ref{thm:tautological} and the measured multipliers (which
inherit the CLARABEL precision).

Extrapolating $r_C(N) \to 0$ as $N \to \infty$ (which holds
empirically with $R^2 \ge 0.78$ in a $1/N$ fit; cf.\
Section~\ref{sec:phase-transition}), the asymptotic framework ceiling is
\[
C_\infty \;:=\; \mu_{\mathrm{LB}} + \lim_{N \to \infty} r_B(N \mid 20)
\;\approx\; 0.380558.
\]
\end{theorem}

\paragraph{Open-gap decomposition.}
With the current LB $0.380299$ and the asymptotic framework ceiling
$C_\infty \approx 0.380558$, the open interval $[0.380299, 0.380871]$
splits as
\[
\underbrace{[0.380299,\; 0.380558]}_{\text{framework-attainable},\;\sim 45\%}
\;\cup\;
\underbrace{[0.380558,\; 0.380871]}_{\text{beyond framework},\;\sim 55\%}.
\]

The framework-attainable portion can in principle be closed by
augmenting the current constraint families and/or scaling $N$; the
beyond-framework portion provably cannot be closed without
fundamentally new techniques.

\section{Phase-transition diagnosis and reciprocal scaling}\label{sec:phase-transition}

The dual multiplier sums $\sum_m m\, \lambda_m^{\cos}(N)$ are not
monotone in $N$ for individual centers: each row has a critical
``activation scale'' $N^*$ at which the lag-$1$ multiplier $\lambda_1$
jumps from approximately zero to approximately $3$--$4$. For three of
our four representative centers this happens at $N^* \approx 17{,}000$;
for the row~7 center (with the largest $h = 0.030$) it is delayed to
$N^* \approx 35{,}000$.

The mechanism is straightforward: $\alpha_1^-(j)$ is a crude
approximation of $\cos(\pi x / 2)$ at small $N$, so the $m=1$ cosine
cell-envelope constraint is slack and $\lambda_1 = 0$. As $N \to \infty$
the constraint tightens and $\lambda_1$ activates. After activation,
multipliers stabilize and the residual decays as $1/N$. This justifies
the extrapolation $r_C(N) \to 0$ used in Theorem~\ref{thm:saturation-formal}.

\section{Numerical results and verification}

Table~\ref{tab:results} summarizes the per-row Phase-5 dual objectives
at $N = 20{,}000$.

\begin{table}[h]
\centering
\small
\begin{tabular}{cccccc}
\toprule
row & center $(h, p)$ & $V_r^{\,\text{center}}$ & semi-$h$ & semi-$p$ & status \\
\midrule
row~1 & $(0.015, 0.381)$ & $0.380255$ & $0.220$ & $0.211$ & non-binding \\
row~4 & $(0.004, 0.3875)$ & $0.380090$ & $0.089$ & $0.083$ & binding regime \\
row~7 & $(0.030, 0.375)$ & $0.381586$ & $0.600$ & $0.751$ & non-binding \\
cde\_n30\_iter3 & $(0.000, 0.3901)$ & $0.380318$ & $0.083$ & $0.075$ & \textbf{binding witness} \\
\bottomrule
\end{tabular}
\caption{Per-row dual objectives at $N=20{,}000$ for four representative
centers (4 of 12). The binding witness is the CDE-discovered
\texttt{cde\_n30\_iter3} center.}
\label{tab:results}
\end{table}

\paragraph{Verification.}
Three independently-written code paths (\texttt{path\_b\_analytical.py},
\texttt{path\_b\_rigorous.py}, \texttt{path\_b\_independent.py}) agree
on per-row SDP solves to $10$+ digits. A separate agent reimplemented
the Bochner-PSD encoding from scratch and obtained bit-for-bit
agreement on test inputs. CLARABEL's primal-dual gap on these solves
is $3 \times 10^{-6}$ to $4 \times 10^{-8}$, comfortably within the
$10^{-6}$ safety margin. Cross-solver agreement between CLARABEL and
SCS at moderate $N$ is $\sim 5 \times 10^{-6}$.

\paragraph{Reproducibility.}
Code, raw SDP outputs, and three independent implementations of the
ellipse-extension argument are in the GitHub repository at
\url{https://github.com/USERNAME/erdos-minimum-overlap-saturation}. The
single-file reproducer for the binding row is
\texttt{lp\_research\_state/code/path\_b\_with\_polymoment.py}.

\section{Open questions and future directions}\label{sec:open}

The saturation theorem (Theorems~\ref{thm:saturation-informal},
\ref{thm:saturation-formal}) establishes that approximately $55\%$ of
the open gap $[0.380299, 0.380871]$ is beyond the reach of the present
framework. Closing this remaining gap requires techniques outside the
constraint families F1--F5. We highlight several candidate directions.

\paragraph{Question 1 (strict complementarity).}
Theorem~\ref{thm:tautological} is a tautological identity, not the
strict-form complementarity inequality. The empirical evidence is
suggestive: at $n = 20 \to n = 30$, cell-envelope multipliers shrink
by $40$--$45\%$ on three of our four representative rows, suggesting
the residuals are coupled. Does the strict-form
\[
r_{CB}(\infty \mid n) \;\le\; \max\bigl(r_C(\infty),\; r_B(\infty \mid n)\bigr)
\]
admit a clean proof in White's framework?

\paragraph{Question 2 (direct sup\_t representation).}
The cell-envelope constraints are the proximate bottleneck. Is there a
cleaner SDP formulation of $\sup_t (h \star h)(t)$ that avoids the
per-$m$ cell envelope altogether? A moment-matrix representation of
$\hat M(k)$ with constraints derived from $M = f \star f$ would be a
natural candidate.

\paragraph{Question 3 (Lasserre hierarchy at level 3).}
Lasserre level 2 was unsuccessful (the natural rigorization via
Fej\'er--Riesz tail bounds offsets the empirical gain at currently
tractable $T_{\max}$). Level 3 has not been explored; the additional
flexibility might rescue a portion of the heuristic gain.

\paragraph{Question 4 (alternative bases).}
Wavelets, Chebyshev polynomials, and Walsh functions all provide
alternative discretizations of $f$ that suppress the Gibbs phenomenon
near $f$'s discontinuities. We have not systematically tested whether
any provides a tighter SDP relaxation than the standard Fourier basis.

\paragraph{Question 5 (Together AI collaboration on the LB side).}
The Together AI / Einstein Arena team \cite{TogetherComputer2026,
TTTDiscover2026} maintains the current state-of-the-art upper bound
via continued numerical optimization of TTT-Discover's 600-piece
construction. The natural complement on the lower-bound side is a
joint research effort with their methodology applied to dual
certificates.

\section{Acknowledgments}

I am grateful to E.~P.~White for the mathematical foundation; this
work would not exist without his Section~5 framework. The
computational and drafting workflow was assisted throughout by
Anthropic's Claude in a multi-agent configuration.

\bibliographystyle{plain}
\begin{thebibliography}{9}

\bibitem{Erdos1955}
P.~Erd\H{o}s.
``Some remarks on number theory'' (in Hebrew).
\textit{Riveon Lematematika}, 9, 45–48, 1955.

\bibitem{Haugland2016}
J.~K.~Haugland.
``A new upper bound on the constant in the Erd\H{o}s minimum overlap problem.''
arXiv:1609.08000, 2016.

\bibitem{White2023}
E.~P.~White.
``A new bound for Erd\H{o}s' minimum overlap problem.''
\textit{Acta Arithmetica} 208 (2023), 235--255. arXiv:2201.05704.

\bibitem{AlphaEvolve2025}
A.~Novikov et al. (Google DeepMind).
``AlphaEvolve: A coding agent for scientific and algorithmic discovery.''
arXiv:2506.13131, May 2025.

\bibitem{TTTDiscover2026}
M.~Yuksekgonul, D.~Koceja, X.~Li, F.~Bianchi, J.~McCaleb, X.~Wang,
J.~Kautz, Y.~Choi, J.~Zou, C.~Guestrin, Y.~Sun.
``Learning to discover at test time.''
arXiv:2601.16175, January 2026.

\bibitem{TogetherComputer2026}
Together AI / Einstein Arena.
``New State-of-the-Art on Erd\H{o}s' Minimum Overlap Problem.''
\url{https://github.com/togethercomputer/erdos-minimum-overlap}, March 2026.

\bibitem{SDPA}
M.~Yamashita, K.~Fujisawa, M.~Kojima.
``SDPARA: SemiDefinite Programming Algorithm paRAllel version.''
\textit{Parallel Computing} 29 (2003), 1053--1067.

\bibitem{LasserreTail}
B.~Zanghi.
``Fej\'er-Riesz tail bound for Lasserre level-2 truncation
in the Erd\H{o}s minimum-overlap SDP.''
Repository note, May 2026.

\end{thebibliography}

\end{document}
